\begin{lemma}
Let $k\ge 3$, and assign a nonnegative real weight $w_{ij}$ to each unordered
pair $\{i,j\}\subseteq \{1,\ldots,k\}$.  Put
\[
M=\sum_{1\le i<j\le k}w_{ij}.
\]
Then
\[
\sum_{1\le i<j<\ell\le k}\sqrt{w_{ij}w_{i\ell}w_{j\ell}}
\le
\binom{k}{3}\left(\frac{M}{\binom{k}{2}}\right)^{3/2}.
\]
\end{lemma}
\begin{proof}
It is enough to prove the claim for $M>0$, since if $M=0$ then every
nonnegative weight is zero and both sides are zero.  Assume $M>0$.
For a vertex $i$ of the index set, write
\[
M_i=\sum_{j\ne i}w_{ij}.
\]
Also write $T_i$ for the part of the weighted triangle sum coming from
triangles that contain $i$:
\[
T_i=\sum_{\substack{j<\ell\\ j,\ell\ne i}}
\sqrt{w_{ij}w_{i\ell}w_{j\ell}}.
\]
Every weighted triangle has exactly three vertices, so
\[
3\sum_{i<j<\ell}\sqrt{w_{ij}w_{i\ell}w_{j\ell}}
=\sum_i T_i. \tag{1}
\]

Fix $i$.  Apply Cauchy--Schwarz to the finite family indexed by unordered
pairs $\{j,\ell\}$ disjoint from $\{i\}$, with entries
$\sqrt{w_{ij}w_{i\ell}}$ and $\sqrt{w_{j\ell}}$.  This gives
\[
T_i^2
\le
\left(\sum_{\substack{j<\ell\\ j,\ell\ne i}} w_{ij}w_{i\ell}\right)
\left(\sum_{\substack{j<\ell\\ j,\ell\ne i}} w_{j\ell}\right). \tag{2}
\]
The second factor in (2) is the total weight of pairs not incident with
$i$, namely $M-M_i$.  For the first factor, the numbers
$(w_{ij})_{j\ne i}$ are $k-1$ nonnegative numbers with sum $M_i$.  Hence
\[
\sum_{\substack{j<\ell\\ j,\ell\ne i}}w_{ij}w_{i\ell}
=\frac{M_i^2-\sum_{j\ne i}w_{ij}^2}{2}
\le
\frac{M_i^2-M_i^2/(k-1)}2
=\frac{k-2}{2(k-1)}M_i^2, \tag{3}
\]
where the inequality uses Cauchy--Schwarz in the form
$\sum_{j\ne i}w_{ij}^2\ge M_i^2/(k-1)$.  Combining (2) and (3), and using
nonnegativity to take square roots, gives
\[
T_i\le
M_i\sqrt{\frac{k-2}{2(k-1)}(M-M_i)}. \tag{4}
\]

Set $x_i=M_i/M$.  Since $M>0$, these numbers are nonnegative, and because
each unordered pair contributes to exactly two of the sums $M_i$, we have
\[
\sum_i x_i=2. \tag{5}
\]
Also $M_i\le M$, so $0\le x_i\le 1$.  Summing (4) over $i$ and factoring
out $M^{3/2}$ gives
\[
\sum_i T_i
\le
M^{3/2}\sqrt{\frac{k-2}{2(k-1)}}
\sum_i x_i\sqrt{1-x_i}. \tag{6}
\]
The remaining finite optimization is another direct Cauchy--Schwarz
calculation:
\[
\left(\sum_i x_i\sqrt{1-x_i}\right)^2
\le
\left(\sum_i x_i\right)
\left(\sum_i x_i(1-x_i)\right). \tag{7}
\]
By (5), the first factor in (7) is $2$.  For the second factor,
\[
\sum_i x_i(1-x_i)=2-\sum_i x_i^2
\le
2-\frac{(\sum_i x_i)^2}{k}
=2-\frac4k
=\frac{2(k-2)}k, \tag{8}
\]
again by Cauchy--Schwarz.  Equations (7) and (8) imply
\[
\sum_i x_i\sqrt{1-x_i}\le 2\sqrt{\frac{k-2}{k}}. \tag{9}
\]
Substituting (9) into (6), then dividing by $3$ using (1), yields
\[
\sum_{i<j<\ell}\sqrt{w_{ij}w_{i\ell}w_{j\ell}}
\le
\frac{\sqrt2}{3}\frac{k-2}{\sqrt{k(k-1)}}\,M^{3/2}. \tag{10}
\]
Finally
\[
\frac{\binom{k}{3}}{\binom{k}{2}^{3/2}}
=
\frac{k(k-1)(k-2)/6}{(k(k-1)/2)^{3/2}}
=
\frac{\sqrt2}{3}\frac{k-2}{\sqrt{k(k-1)}}.
\]
Replacing the coefficient in (10) by this identical binomial expression is
exactly the desired bound.
\end{proof}
